% =====================================================================
% BMED365: Explainable AI (XAI)
% Beamer-presentasjon - Topic List X01-X08
% =====================================================================
\documentclass[aspectratio=169, 10pt]{beamer}

% =====================================================================
% PACKAGES
% =====================================================================
\usepackage[utf8]{inputenc}
\usepackage[T1]{fontenc}
\usepackage[english]{babel}
\usepackage{graphicx}
\usepackage{tikz}
\usetikzlibrary{shapes.geometric, arrows, positioning, calc}
\usepackage{booktabs}
\usepackage{amsmath}
\usepackage{fontawesome5}
\usepackage{hyperref}
\hypersetup{
    colorlinks=true,
    linkcolor=blue,
    urlcolor=blue,
    citecolor=blue
}

% =====================================================================
% THEME AND COLORS
% =====================================================================
\usetheme{Madrid}
\usecolortheme{beaver}

% =====================================================================
% TITLE INFO
% =====================================================================
\title{Explainable AI (XAI)}
\subtitle{BMED365: Topic List X01--X08}
\author{BMED365}
\date{Spring 2026}

% =====================================================================
% DOCUMENT
% =====================================================================
\begin{document}

% Title page
\begin{frame}
    \titlepage
\end{frame}

% Table of Contents
\begin{frame}{Overview}
    \begin{enumerate}
        \item \textbf{Why explainability?}
        \begin{itemize}
            \item X01: Explain why explainability is important in medical AI
        \end{itemize}
        \item \textbf{Types of Explainability}
        \begin{itemize}
            \item X02: Distinguish between global and local explainability
            \item X03: Distinguish between ante-hoc and post-hoc explainability
        \end{itemize}
        \item \textbf{XAI-metoder}
        \begin{itemize}
            \item X04: Describe SHAP (SHapley Additive exPlanations)
            \item X05: Describe LIME (Local Interpretable Model-agnostic Explanations)
            \item X06: Explain Grad-CAM for CNN visualization
        \end{itemize}
        \item \textbf{Limitations and Challenges}
        \begin{itemize}
            \item X07: Discuss limitations of XAI methods
            \item X08: Know about attention visualization in LLMs
        \end{itemize}
    \end{enumerate}
\end{frame}

% =====================================================================
% SECTION: Why explainability?
% =====================================================================
\section{Why explainability?}

\begin{frame}{X01: Explain why explainability is important in medical AI}
    \textbf{Challengeen med ``\href{https://en.wikipedia.org/wiki/Black_box}{black box}'' AI:}
    \begin{itemize}
        \item Komplekse models (dype nettverk, LLM) gir predictioner without explain \textit{hvorfor}
        \item In medicine er dette problematisk av flere grunner
    \end{itemize}
    
    \vspace{0.2cm}
    \textbf{Why explainability er critical i medisin:}
    \begin{enumerate}
        \item \textbf{Tillit:} Leger og patienter må stole på anbefalingene
        \item \textbf{Responsibility:} Ved feil -- hvem har the responsibility?
        \item \textbf{Regulation:} EU AI Act krever explainability for høyrisiko-AI
        \item \textbf{Læring:} Understand hva modellen har lært (og feil-lært)
        \item \textbf{Feilsøking:} Identifisere bias og weaknesser
        \item \textbf{Clinical integration:} AI som beslutningsstøtte, ikke autonom
    \end{enumerate}
    
    {\footnotesize
    \begin{alertblock}{Medical practice}
        En lege som ikke kan explain why de anbefaler en treatment vil miste patientens tillit. Det samme gjelder AI.
    \end{alertblock}
    }
\end{frame}

% =====================================================================
% SECTION: Types of Explainability
% =====================================================================
\section{Types of Explainability}

\begin{frame}{X02: Distinguish between global and local explainability}
    \begin{columns}[T]
        \begin{column}{0.48\textwidth}
            \textbf{Global explainability:}
            \begin{itemize}
                \item Explains modellens \textbf{generelle behavior}
                \item ``Hvilke features er importantst totalt sett?''
                \item ``How påvirker alder predictionen generelt?''
            \end{itemize}
            
            \vspace{0.2cm}
            \textbf{Metoder:}
            \begin{itemize}
                \item \href{https://en.wikipedia.org/wiki/Feature_importance}{Feature importance}
                \item \href{https://christophm.github.io/interpretable-ml-book/pdp.html}{Partial dependence plots}
                \item \href{https://shap.readthedocs.io/en/latest/example_notebooks/api_examples/plots/beeswarm.html}{Globale SHAP-valueer}
            \end{itemize}
        \end{column}
        \begin{column}{0.48\textwidth}
            \textbf{Lokal explainability:}
            \begin{itemize}
                \item Explains én \textbf{spesifikk prediction}
                \item ``Why fikk \textit{denne} patienten høy risikoscore?''
                \item ``Hvilke piksler påvirket classificationen?''
            \end{itemize}
            
            \vspace{0.2cm}
            \textbf{Metoder:}
            \begin{itemize}
                \item \href{https://github.com/marcotcr/lime}{LIME}
                \item \href{https://shap.readthedocs.io/}{SHAP} (individuelle valueer)
                \item \href{https://arxiv.org/abs/1610.02391}{Grad-CAM} (for bilder)
            \end{itemize}
        \end{column}
    \end{columns}
    
    \vspace{0.2cm}
    {\small
    \begin{block}{Begge er importante}
        \textbf{Global:} Understand modellen som helhet, validere mot domenkunnskap \\
        \textbf{Lokal:} Explain beslutninger til patienter og kollegaer
    \end{block}
    }
\end{frame}

\begin{frame}{X03: Distinguish between ante-hoc and post-hoc explainability}
    \begin{columns}[T]
        \begin{column}{0.48\textwidth}
            \textbf{\href{https://christophm.github.io/interpretable-ml-book/simple.html}{Ante-hoc} (innebygd):}
            \begin{itemize}
                \item Explainability \textbf{bygget inn} i modellen
                \item Modellen er designet to være tolkbar
                \item ``Interpretable by design''
            \end{itemize}
            
            \vspace{0.2cm}
            \textbf{Examples:}
            \begin{itemize}
                \item Lineær regression (koeffisienter)
                \item Beslutningstrær (regler)
                \item Attention-vekter (delvis)
            \end{itemize}
            
            \vspace{0.2cm}
            \textbf{Fordel:} Ekte forklaring \\
            \textbf{Ulempe:} Ofte mindre nøyaktig
        \end{column}
        \begin{column}{0.48\textwidth}
            \textbf{\href{https://christophm.github.io/interpretable-ml-book/agnostic.html}{Post-hoc} (etterpå):}
            \begin{itemize}
                \item Forklaringsmetode \textbf{layerst til} etter training
                \item Modellen er fortsatt en ``black box''
                \item Separate verktøy for interpretation
            \end{itemize}
            
            \vspace{0.2cm}
            \textbf{Examples:}
            \begin{itemize}
                \item SHAP, LIME
                \item Grad-CAM
                \item \href{https://en.wikipedia.org/wiki/Saliency_map}{Saliency maps}
            \end{itemize}
            
            \vspace{0.2cm}
            \textbf{Fordel:} Kan brukes på complexe models \\
            \textbf{Ulempe:} Approksimerer, can be misvisende
        \end{column}
    \end{columns}
    
    \vspace{0.1cm}
    {\small
    \begin{alertblock}{Trade-off}
        Ofte et valg mellom tolkbarhet (ante-hoc) og ytelse (complex modell + post-hoc)
    \end{alertblock}
    }
\end{frame}

% =====================================================================
% SECTION: XAI-metoder
% =====================================================================
\section{XAI-metoder}

\begin{frame}{X04: Describe SHAP (SHapley Additive exPlanations)}
    \textbf{SHAP: Basert på spillteori}
    
    \begin{itemize}
        \item Shapley-valueer: Fair advantageing av ``gevinst'' mellom spillere
        \item Her: Hvor mye bidrar hver feature til predictionen?
    \end{itemize}
    
    \vspace{0.2cm}
    \textbf{How SHAP works:}
    \begin{enumerate}
        \item For hver feature: Beregn gjennomsnittlig bidrag til prediction
        \item Tar hensyn til alle mulige kombinasjoner av features
        \item Gir en value per feature for hver prediction
    \end{enumerate}
    
    \vspace{0.15cm}
    \begin{columns}[T]
        \begin{column}{0.48\textwidth}
            \textbf{Styrker:}
            \begin{itemize}
                \item Teoretisk solid (\href{https://en.wikipedia.org/wiki/Shapley_value}{Shapley-aksiomer})
                \item Konsistent og lokalt nøyaktig
                \item Fungerer for alle modelltyper
            \end{itemize}
        \end{column}
        \begin{column}{0.48\textwidth}
            \textbf{Visualiseringer:}
            \begin{itemize}
                \item \textbf{\href{https://shap.readthedocs.io/en/latest/example_notebooks/api_examples/plots/waterfall.html}{Waterfall plot}:} Én prediction
                \item \textbf{\href{https://shap.readthedocs.io/en/latest/example_notebooks/api_examples/plots/beeswarm.html}{Summary plot}:} Global oversikt
                \item \textbf{\href{https://shap.readthedocs.io/en/latest/example_notebooks/api_examples/plots/force.html}{Force plot}:} Interaktiv
            \end{itemize}
        \end{column}
    \end{columns}
    
    {\footnotesize
    \begin{block}{Python}
        \texttt{import shap; explainer = shap.Explainer(model); shap\_values = explainer(X)}
    \end{block}
    }
\end{frame}

\begin{frame}{X05: Describe LIME (Local Interpretable Model-agnostic Explanations)}
    \textbf{\href{https://github.com/marcotcr/lime}{LIME}: Lokal tilnærming med simple modell}
    
    \vspace{0.2cm}
    \textbf{Main idea:}
    \begin{enumerate}
        \item Velg én prediction å explain
        \item Generer \textbf{\href{https://en.wikipedia.org/wiki/Perturbation_theory}{perturberte}} (fotherwisetyrrede) versjoner av input
        \item Kjør den complexe modellen på alle versjoner
        \item Tren en \textbf{simple, tolkbar modell} (e.g. lineær) lokalt
        \item Bruk den enkle modellen til å explain
    \end{enumerate}
    
    \vspace{0.2cm}
    \begin{columns}[T]
        \begin{column}{0.48\textwidth}
            \textbf{Styrker:}
            \begin{itemize}
                \item Modell-agnostisk
                \item Intuitivt understandelig
                \item Fungerer for tekst, bilder, tables
            \end{itemize}
        \end{column}
        \begin{column}{0.48\textwidth}
            \textbf{Svakheter:}
            \begin{itemize}
                \item Ustabil -- ulike kjøringer kan gi ulike svar
                \item Valg av nabolayers er vilkårlig
                \item Kan være misvisende
            \end{itemize}
        \end{column}
    \end{columns}
    
    {\footnotesize
    \begin{block}{Example: Tekstclassification}
        LIME viser which ord som bidro mest til classificationen ``spam'' vs. ``ikke spam''
    \end{block}
    }
\end{frame}

\begin{frame}{X06: Explain Grad-CAM for CNN visualization}
    \textbf{\href{https://arxiv.org/abs/1610.02391}{Grad-CAM}: Gradient-weighted Class Activation Mapping}
    
    \begin{itemize}
        \item Spesialisert for \textbf{konvolusjonelle nettverk (CNN)}
        \item Viser \textbf{which bildområder} som påvirket classificationen
    \end{itemize}
    
    \vspace{0.2cm}
    \textbf{How it works:}
    \begin{enumerate}
        \item Beregn gradienter av output with respect to feature maps i siste konvolusjonslayers
        \item Vekt feature maps med gjennomsnittlig gradient
        \item Kombiner til et \textbf{\href{https://en.wikipedia.org/wiki/Heat_map}{heatmap}} som overlays på originalbildet
    \end{enumerate}
    
    \vspace{0.2cm}
    \textbf{Medical application:}
    \begin{itemize}
        \item Røntgen/CT: ``Modellen fokuserte på dette området for lungefunn''
        \item Dermatologi: ``Disse pikslene påvirket melanom-diagnosisn''
        \item Validation: Sjekk at modellen ser på riktig område!
    \end{itemize}
    
    {\footnotesize
    \begin{alertblock}{Important innsikt}
        Grad-CAM kan avsløre om modellen bruker \textbf{spuriøse korrelasjoner} (e.g. ser på tekst i bildet, ikke patologi)
    \end{alertblock}
    }
\end{frame}

% =====================================================================
% SECTION: Limitations and Challenges
% =====================================================================
\section{Limitations and Challenges}

\begin{frame}{X07: Discuss limitations of XAI methods}
    \textbf{XAI er ikke perfekt -- \href{https://arxiv.org/abs/2006.11371}{importante limitationer}:}
    
    \vspace{0.2cm}
    \begin{enumerate}
        \item \textbf{Forklaringer er approksimeringer}
        \begin{itemize}
            \item Post-hoc metoder explains ikke den \textit{ekte} modellen
            \item Kan være misvisende or inkonsistente
        \end{itemize}
        
        \item \textbf{Ustabilitet}
        \begin{itemize}
            \item LIME kan gi ulike forklaringer for samme input
            \item Små endringer i input kan gi scountse endringer i forklaring
        \end{itemize}
        
        \item \textbf{Forinterpretationsbyrde}
        \begin{itemize}
            \item Hvem shall interpret forklaringene? Krever ekspertise
            \item Risiko for overtillit til forklaringer
        \end{itemize}
        
        \item \textbf{Beregningsmessig kost}
        \begin{itemize}
            \item SHAP can be svært treg for scountse models
        \end{itemize}
    \end{enumerate}
    
    \vspace{0.1cm}
    {\small
    \begin{block}{Husk}
        XAI er et verktøy for innsikt, ikke en garanti for at modellen er ``riktig'' or rettferdig.
    \end{block}
    }
\end{frame}

\begin{frame}{X08: Know about attention visualization in LLMs}
    \textbf{\href{https://jalammar.github.io/illustrated-transformer/}{Attention-visualisering} i LLM: Innebygd ``forklaring''?}
    \vspace{-0.1cm}
    \begin{itemize}
        \item Transformer-models bruker \textbf{\href{https://en.wikipedia.org/wiki/Attention_(machine_learning)}{self-attention}}
        \item Attention-vekter viser which tokens modellen ``ser på''
        \item Kan visualiseres som heatmaps
    \end{itemize}
    
    \vspace{0.1cm}
    \textbf{Example:}
    \begin{center}
        \textit{``Pasienten har \textbf{diabetes} og tar \textbf{metformin}''}
    \end{center}
    Attention-visualisering kan vise at ``metformin'' har høy attention på ``diabetes''
    
    \vspace{0.1cm}
    \begin{columns}[T]
        \begin{column}{0.48\textwidth}
            \textbf{Muligheter:}
            \begin{itemize}
                \item Innsikt i modellens fokus
                \item Ante-hoc (innebygd i modellen)
                \item Lett tilgjengelig
            \end{itemize}
        \end{column}
        \begin{column}{0.48\textwidth}
            \textbf{Limitations:}
            \begin{itemize}
                \item Attention $\neq$ forklaring
                \item Mange layers, mange \href{https://en.wikipedia.org/wiki/Transformer_(deep_learning_architecture)}{``heads''}
                \item Kan være vanskelig å interpret
                \item Viser korrelasjon, ikke kausalitet
            \end{itemize}
        \end{column}
    \end{columns}
    
    {\footnotesize
    \begin{alertblock}{Fotherwiseiktighet}
        Attention-visualiseringer bør brukes som supplement, ikke som endelig forklaring.
    \end{alertblock}
    }
\end{frame}

% =====================================================================
% SUMMARY
% =====================================================================
\section*{Summary}

\begin{frame}{Summary: Explainable AI (XAI)}
    \textbf{Key points:}
    \begin{itemize}
        \item \textbf{X01:} Explainability er critical i medisin (tillit, responsibility, regulation)
        \item \textbf{X02:} Global vs. lokal explainability
        \item \textbf{X03:} Ante-hoc (innebygd) vs. post-hoc (etterpå)
        \item \textbf{X04:} SHAP -- spillteoretisk, konsistent
        \item \textbf{X05:} LIME -- lokal tilnærming, modell-agnostisk
        \item \textbf{X06:} Grad-CAM -- visualiser CNN-fokus i bilder
        \item \textbf{X07:} Limitations -- approksimeringer, ustabilitet
        \item \textbf{X08:} Attention-visualisering i LLM
    \end{itemize}
    
    \vspace{0.2cm}
    {\small
    \begin{block}{Practical work in Lab 2 og Lab 3}
        \begin{itemize}
            \item Lab 2: Bruk Grad-CAM to interpret CNN-classification
            \item Lab 3: Diskuter explainability og transparens i LLM
        \end{itemize}
    \end{block}
    }
\end{frame}

\end{document}






